\section{Audit, integration corrections, and further research}
\label{audit:section}

\subsection{What the audit changes}

No new false theorem in the focused established source material was
identified. The principal audit finding is a status issue: several
questions survive in earlier archives after later incoming reports have
answered them. Integration should reconcile the chronology before
assigning credit or declaring a question newly resolved.

\begin{longtable}{>{\raggedright\arraybackslash}p{0.29\textwidth}
                  >{\raggedright\arraybackslash}p{0.63\textwidth}}
\toprule
Earlier question or claim & Proposed integration treatment\\
\midrule
\endhead
Single-resolvent nonpositive powers &
The complete answer is already in \cite{repo:quartic}. The earlier
Barnes Question F5 is therefore resolved by that report. The new
theorem here concerns independently powered products and their
orientation-dependent endpoint poles.\\[5pt]
The fourth digamma power &
Retain the existing cubic and quartic bridges with credit to
\cite{repo:cubic,repo:barnes,repo:quartic}. Add the fifth theorem and the
all-power reduction as subsequent extensions.\\[5pt]
Pure pairs of even harmonic tails &
The question in the earlier parity and reflection reports is already
answered by \cite{repo:quartic}. It is not a new contribution of this
article. Products with further factors remain a separate problem.\\[5pt]
Higher-index collision remainder &
Replace the open-status entry for existence of a fully subtracted
total limit by Theorem~\ref{gcx:main}. State continuity explicitly.
Do not upgrade this to analyticity at positive Stieltjes indices or
to associativity of arbitrary partial-merger prescriptions.\\[5pt]
Canonical power polynomials &
Record Corollary~\ref{hp:cancellation} as elimination of the inherited
height-one coefficients $\eta,\theta$. The additional derivative
coordinate $\chi_4$ remains, so the stronger single-coordinate
all-power question is still open.\\[5pt]
Gaussian $S_6$ and revised $S_8$ &
Preserve their conjectural status. Keep the previously rejected
coefficient vector distinct from the revised candidate, and distinguish
period equality from membership in a specified formal relation system.\\
\bottomrule
\end{longtable}

The companion \src{integration/CLAIM_LEDGER.csv} records theorem labels,
source questions, and proposed status changes. The source snapshot
manifest gives the inspected revision and hashes of the current incoming
archives. The older reports are dependencies and provenance, not copied
into this deliverable as if they were new work.

\subsection{Exact corrections to tempting extrapolations}

Two natural extrapolations are disproved by the calculations here.
They should be excluded explicitly in any consolidated statement.

\paragraph{A value along a parameter path is not automatically the
coordinate finite part.}
For opposite resolvents, \eqref{mr:pathlimit} gives the three limits
$-2\pi\ii,-\pi\ii,0$ for the same specialized function $K-\pi\ii$.
Its coordinate finite part is $-\pi\ii$.
The proposed replacement is the full germ \eqref{mr:H}, retaining
$E=\sum\epsilon_j\lambda_j$ as well as $\Lambda=\sum\lambda_j$.
The phase numerator may be simplified using integrality of each exponent
only when that hypothesis actually holds. The theorem itself requires
only integrality of the total at the chosen resonance.

\paragraph{Taking the order-zero term of a collision does not commute
with merging its factors.}
The three constants \eqref{gcx:CT0}--\eqref{gcx:CT2} differ although
their directly merged integrand is the same
$\gamma_1\gamma_0^2$. The replacement is to subtract the entire
counterterm \eqref{gcx:T} before passing to the limit. Dividing the local
cutoff by a shrinking gap would alter this convention and its constants.
The terms involving
$\gamma_1'(1)=\zeta(2)+\zeta'(2)$ must retain both summands;
differentiation of the Pochhammer factor supplies the first one.

These counterexamples correct unsupported extensions. They are not
retroactive claims that the earlier single-resolvent or index-zero
theorems were false.

Two further proof requirements should accompany the integration.
Matching principal parts of $\psi^k$ leaves an entire-function
ambiguity; Lemma~\ref{hp:ML} supplies the contour argument that eliminates
it. Also, an exact finite harmonic identity is not the same statement as
its asymptotic regularization: use \eqref{hp:regularization} with finite
terminal harmonic sums first, and replace convergent terminal words by
their limits only after controlling their tails. These requirements are
met in the proofs above.

\subsection{Questions opened by the new results}

The following projects ask for exact identities or structural theorems.
They are questions rather than unproved assertions of reducibility.
Each has a concrete starting formula in this article.

\begin{question}[The derivative coordinate of weight four]
Can
\[
 \chi_4=2(\partial_1-\partial_2)Z_2(3,1)+\partial_1Z_2(2,2)
\]
be reduced to a smaller specified algebra of ordinary zeta derivatives
and height-one Laurent coefficients?
\end{question}
The first task is to differentiate identities valid with genuinely free
orders, rather than differentiate a duality known only at discrete
integer indices. The free-order stuffle identity used after
\eqref{hp:Er} is admissible in this sense. A useful result would be an
exact reduction or a formal separating functional for a precisely
enumerated family of differentiated relations. Such a formal obstruction
would still not prove arithmetic independence.

\begin{question}[The sixth power and canonical cancellation matrices]
Compute the complete sixth-power reduction and determine which
polynomials in $\psi$ eliminate its inherited lower-depth coordinates.
Does a uniform family of power polynomials remove all but the single
height-one coefficient at every degree?
\end{question}
The compiler of Theorem~\ref{hp:allpowers} already terminates at every
fixed degree and fixes its constants. What remains is to simplify the
resulting moving-order coordinates and compute the exact cancellation
matrix against lower powers. The coefficient of the top-depth term is
already known, namely $-k!$; it is the lower-depth relations that must be
understood. The surviving $\chi_4$ in the fifth case is a specific test.

\begin{question}[A common Hurwitz shift for all digamma powers]
Extend the reduction and the normalized primitive to
$\psi(x+a)^k$, with a free common shift $a$, and determine the finite
multiple-Hurwitz or Lerch family required.
\end{question}
At $a>0$ the unit-interval integral is ordinary; the principal poles of
the global meromorphic function have moved to $-a-\N_0$.
The contour subtraction and finite summation-by-parts method should be
carried out with the shifted lattice, keeping the endpoint normalization
explicit. Rational $a$ offers a precise next specialization to colored
multiple-zeta data.

\begin{question}[Higher mixed logarithmic resolvent products]
For arbitrary positive integers $r,s$, find finite multiple-polylogarithmic
expressions for the all-index transforms of $L_z^rL_w^s$.
What depth and moving-order derivative families are sufficient?
\end{question}
The two-factor formula \eqref{ml:Tu} and the ordinary cubic formula
\eqref{ml:triple} are initial cases. Start from the integer convolution
constraints of the Fourier coefficients, apply exact partial fractions,
and keep the complete local power--log polynomial in the Stieltjes
subtraction. A finite representation in a stated class is the target;
introducing an unrestricted new generating integral alone would not
provide that reduction.

\begin{question}[Classical multivariable forms of the Gamma spectrum]
For which exponent patterns and disk coordinates does
\eqref{mr:ck} collapse to a finite combination of Gauss, Appell,
Lauricella, or multiple-polylogarithmic functions?
\end{question}
The example \eqref{mr:gaussmean} proves that a Gauss function suffices
for one zero disk coordinate. The general Gamma convolution is already
explicit and normally convergent. A further result should produce a
smaller function class, a new functional equation, or an arithmetic
special-value identity, with all exceptional parameter restrictions
handled on the complete meromorphic expression.

\begin{question}[Real-axis boundary identities]
Determine the exact limits of the completed mixed resolvent transforms
as one or more nonreal parameters approach the real axis, including
coincident limiting parameters.
\end{question}
The sign $\epsilon_j$ then records which boundary value is taken, and
interior cotangent poles can appear. A correct statement must inventory
their local principal-value and contact contributions as well as the
endpoint terms already present in \eqref{mr:H}. This is an exact
boundary-value problem, not a request for numerical bounds.

\begin{question}[Regularity of the higher-index collision remainder]
Find the complete power--log expansion of $C-T$ near total collision,
and determine its differentiability and analytic subfamilies.
\end{question}
Theorem~\ref{gcx:main} proves continuity uniformly in the relative
gaps. It does not assert a convergent Taylor series. After the Hermite
subtraction the only obstacle is the explicit product of translated
logarithms. Resolving that remaining integral for one cluster and then
for nested clusters would give a controlled higher-order contact
calculus, including gap configurations not described by a single scale.

\begin{question}[The global mixed cubic constant]
Express
$Q=\FP\int_0^1\gamma_1(x)\gamma_0(x)^2\,\dd x$
in a specified multiple-Hurwitz or directional Tornheim derivative
family, and seek combinations that eliminate it.
\end{question}
Equation~\eqref{gcx:Q-ordinary} is an absolutely convergent representative.
The difference of collision constants already eliminates $Q$ exactly.
For a further arithmetic reduction, begin with a joint completed
three-factor Hurwitz generator and differentiate its orders before
specialization. A numerical relation for the displayed decimal would
only suggest a candidate.

\begin{question}[Compatibility of successive partial mergers]
Characterize partial-cluster subtraction rules whose iterated mergers
agree with the simultaneous full counterterm \eqref{gcx:T}.
Can a universal local adjustment make the operation associative?
\end{question}
The present result fixes the total limit even for hierarchical gaps.
It does not identify the finite constants produced by every possible
intermediate prescription. For three and four factors, compare two
merger orders symbolically using the explicit dilogarithm kernels,
then determine the local difference before proposing a general rule.

\begin{question}[Integer-frequency grids and their logarithmic scales]
Extend the geometric theorem to
$\prod_i\gamma_{m_i}^{(r_i)}(\{p_i x+s_i\})$ with positive integers
$p_i$, including analytic Fourier weights.
\end{question}
The analytic-weight portion at common unit frequency is proved here.
For a grid point with ambient displacement $t$, the singular factor
has amplitude
$p_i^{-r_i-1}t^{-r_i-1}L_{m_i,r_i}(\log t+\log p_i)$.
A complete theorem should inventory all coinciding grid points and
transport the cutoff in this ambient coordinate. The logarithms
$\log p_i$ are finite contacts that cannot be removed by an informal
rescaling.

\begin{question}[A common spectral and collision completion]
Construct a joint generator with both the complex resolvent-power
weight and moving periodic Stieltjes factors, and prove compatible
restriction formulas at intersecting spectral and geometric resonances.
\end{question}
The two ingredients are explicit here, but their direct combination
has new local products: a weighted collision point can approach an
endpoint where the weight is singular. The target is one full
subtraction rule that recovers \eqref{mr:H} and \eqref{gcx:T} in
their separate domains, with the order of all restrictions specified.

\begin{question}[A new exact relation for the frozen Gaussian targets]
Can an identity obtained from a harmonic, Stieltjes, or resolvent
transform supply a genuine functional relation for the frozen
$S_6$ or revised $S_8$ residual outside the previously obstructed
formal systems?
\end{question}
The source vectors and branch conventions must remain fixed. A useful
candidate should have a demonstrably new action on the existing
separating functional, or a proved analytic specialization beyond the
enumerated relation family. Neither additional digits nor an
uncompleted modular elimination establishes the period equality.
No such functional bridge has been proved in this article.

\begin{question}[Reusable formalization units]
Formalize the full-amplitude completion lemma, the all-power
Mittag--Leffler normalization, strict harmonic quasi-shuffles and their
asymptotic regularization, and the fixed-contour Hermite remainder.
\end{question}
These units separate the reusable analytic proof from finite symbolic
coefficient calculations. The supplied programs give regression
examples for the latter. Formal verification of the analytic lemmas
would support generated identities at many new indices; importing
Python outputs alone should not confer a proof-assistant status.

\subsection{Suggested order of work}

The most direct next exact calculations are the sixth-power normal
form, the first higher mixed logarithmic transform, and a joint
Hurwitz representation of the global mixed cubic constant. Their input
algebra and all local corrections are now explicit. The collision
regularity and merger-compatibility questions are the next structural
problems. The Gaussian conjectures require a new functional relation or
an exact period certificate beyond the relation systems already
analyzed in the source material.
